\section{Why SETA is the Ceiling: Diagnostic Analysis}
\label{sec:ceiling}

Having presented SETA's recipe and result (Sec.~\ref{sec:recipe}),
we now answer the natural follow-up question: \emph{why does SETA
sit where it does?} In this section we provide two complementary
diagnostic tools --- the \emph{identity-engagement entanglement}
(IDEP) curve and the \emph{within-subject ranking} analysis ---
that quantitatively explain the structural ceiling at
$\kappa_q \approx 0.25$. These diagnostics are themselves
contributions: they are new tools for analysing any feature space
on a subject-disjoint affective recognition benchmark.

\subsection{Experimental setup}
\label{sec:setup}

\textbf{Dataset.} DAiSEE~\cite{gupta2016daisee} provides 8,925
ten-second classroom clips of 112 students annotated with ordinal
engagement labels in $\{0, 1, 2, 3\}$ (\emph{very low, low,
engaged, highly engaged}). We use the official subject-disjoint
splits: 5,358 train (69 subjects), 1,429 validation (19 subjects),
1,784 test (21 subjects). The label distribution is skewed
($L0{:}L1{:}L2{:}L3 = 1\%{:}5\%{:}50\%{:}44\%$ on test). All
results reported on the \emph{full} test set (no sampling).

\textbf{Frames.} We extract three frames per clip at
$t \in \{2, 5, 8\}$~s, encode each independently with the chosen
encoder, and report results for both the single mid-frame ($t=5$~s)
and the temporal mean. Face bounding boxes are detected with
MediaPipe FaceLandmarker at 99.96\% recall; we re-detect blendshapes,
3-d head pose, and 6-d eye gaze for use in
Sec.~\ref{sec:dream}.

\textbf{Encoders.} We probe six encoders spanning two scales and
two pre-training objectives:
\begin{itemize}\itemsep0pt
  \item Small (B/32 or B/14, 512--768~d):
  CLIP-B/32~\cite{radford2021clip}, DINOv2-B/14~\cite{oquab2023dinov2}.
  \item Large (L or SO400M, 768--1152~d): CLIP-L/14,
  SigLIP-L/16~\cite{zhai2023siglip}, SigLIP-SO400M-14, and
  TIPSv2-B/14~\cite{cao2026tipsv2}.
\end{itemize}
SigLIP-L/16 is our reference encoder for the headline; we report
results across all six in ablations (Tab.~\ref{tab:encoders}).

\textbf{Metrics.} We report quadratic-weighted Cohen's $\kappa$
($\kappa_q$, our primary), accuracy, macro-F1, and bootstrap 95\%
confidence intervals (1{,}000 resamples). Per Reviewer~mfAP of our
prior submission, we drop GPT-4o chain-of-thought from headline
tables (98\% refusal rate; numbers preserved with refusal columns
in Tab.~\ref{tab:gpt4o-refusals}).

\textbf{Baselines.} Eight method families are evaluated:
\begin{itemize}\itemsep0pt
  \item \textsc{ZS-Prompt}: zero-shot prompted classification with
  CLIP/LLaVA/GPT-4o at three prompt variants (minimal,
  rubric-anchored, chain-of-thought).
  \item \textsc{LinProbe}: class-balanced logistic regression on
  frozen features, C-tuned on val.
  \item \textsc{OrdProbe}: ridge ordinal regression rounded to
  nearest integer class; CORN ordinal~\cite{shi2023corn}.
  \item \textsc{Bag}: K=20 bootstrap-bag of LR per
  Sec.~\ref{sec:recipe}, with 5 outer seeds.
  \item \textsc{RSB}: random feature-subspace bagging (fraction
  0.40--0.45), inspired by random forests applied to LR.
  \item \textsc{Cal}: content-free image bias subtraction
  \cite{zhao2021calibrate} and class-prior
  adjustment~\cite{menon2021logit}.
  \item \textsc{Contrastive}: SupCon~\cite{khosla2020supcon} and
  subject-stratified InfoNCE projection heads.
  \item \textsc{Adapter}: light learnable bottleneck adapter
  (latent dim $d \in \{64, 128, 256, 512\}$) trained with focal
  ordinal loss.
\end{itemize}

\subsection{Result: every readout-level family hits the same wall}
\label{sec:exhaustion}

Tab.~\ref{tab:exhaustion} reports the best $\kappa_q$ per method
family on SigLIP-L features, including the variants we found in
prior literature (calibrated CLIP~\cite{wortsman2022robust},
SupCon-projection~\cite{khosla2020supcon}, ordinal
CORN~\cite{shi2023corn}). Fig.~\ref{fig:family_strip} shows the
distribution of $\kappa_q$ within each family across our 60+
variants.

\begin{table}[t]
\caption{Best $\kappa_q$ per method family on DAiSEE test
(SigLIP-L; 1,784 clips; bootstrap 95\% CI from 1,000 resamples).
$\Delta$ is the lift over the strongest zero-shot baseline.}
\label{tab:exhaustion}
\centering
\small
\begin{tabular}{lccc}
\toprule
Family & Best variant & $\kappa_q$ & 95\% CI \\
\midrule
\textsc{ZS-Prompt} & CLIP P1 minimal & 0.014 & $[-.014, .041]$\\
\textsc{ZS-Prompt} & LLaVA-1.5-7B P3 CoT & 0.104 & $[.061, .150]$\\
\textsc{ZS-Prompt} & GPT-4o P1 minimal & 0.021 & $[-.014, .060]$\\
\textsc{LinProbe} & SigLIP-L LR (C=1) & 0.199 & $[.156, .242]$\\
\textsc{OrdProbe} & SigLIP-L Ridge ord & 0.179 & $[.140, .216]$\\
\textsc{Bag} & K=20$\times$5seeds LR & 0.213 & $[.163, .259]$\\
\textsc{Bag+Thresh} & K=20$\times$5seeds + tuned $t$ & \textbf{0.238} & $[.191, .284]$\\
\textsc{RSB} & frac=0.45 RSB-LR & 0.214 & $[.168, .257]$\\
\textsc{Cal (CBU)} & CLIP P3 + CBU & 0.104 & $[.061, .150]$\\
\textsc{Contrastive} & SupCon (DINOv2) & 0.131 & $[.087, .172]$\\
\textsc{Adapter} & lat-64, ord-2.0, $t$-tuned & 0.234 & $[.186, .278]$\\
\midrule
\multicolumn{2}{l}{\emph{Naive baseline argmax (SigLIP-L LR):}} & 0.199 & --- \\
\multicolumn{2}{l}{\emph{Ceiling (best across all families):}} & 0.238 & --- \\
\multicolumn{2}{l}{$\Delta$ (best $-$ argmax):} & $+0.039$ & --- \\
\multicolumn{2}{l}{\emph{Supervised SOTA reference (ViBED-Net):}} & $\sim 0.55$ & --- \\
\bottomrule
\end{tabular}
\end{table}

The variance between method families is small relative to the
distance to the supervised SOTA. Sophisticated methods
(SupCon-projection, CORN ordinal regression, adversarial
gradient-reversal heads) do not exceed simple bagged LR with
threshold tuning. Ensemble methods that introduced large gains in
other settings (RSB at fraction 0.45, multi-encoder fusion) move
the needle by $<\!0.01$ once threshold tuning is held constant.

\subsection{Mechanism: identity-engagement entanglement}
\label{sec:idep}

To understand why the readout-level ceiling is so flat, we
quantify how identity and engagement co-locate in the frozen
feature space.

\paragraph{Subject identity is near-perfectly linearly recoverable.}
A class-balanced multinomial logistic regression on within-train
80/20 clip-level splits recovers subject identity at 99.5\% on
SigLIP-L (99.8\% on CLIP-L/14, 99.6\% on DINOv2-B/14). Identity is
not just present in frozen features; it is dominant.

\paragraph{The asymmetric entanglement curve.} For a chosen
encoder, we train a 109-way subject-ID logistic-regression
classifier and rank its weight rows by L2 norm to obtain a basis
of \emph{identity-dominant directions} $\{u_1, u_2, \dots, u_{109}\}$.
For each $k \in \{0, 1, \dots, 69\}$ we project features onto the
orthogonal complement of $\mathrm{span}(u_1, \dots, u_k)$ and
re-evaluate (a)~subject-ID accuracy and (b)~engagement~$\kappa_q$
with a class-balanced LR probe. The result, plotted in
Fig.~\ref{fig:idep}, exhibits an asymmetric structure:

\begin{itemize}\itemsep0pt
\item \textbf{Identity is high-rank redundant.} Subject-ID
accuracy on SigLIP-L is still $99.5\%$ at $k = 5$, $98.5\%$ at
$k = 10$, $97.3\%$ at $k = 20$, and only collapses below 50\% at
$k \geq 50$. The encoder spreads identity across dozens of
near-equivalent directions.
\item \textbf{Engagement is low-rank fragile.} Engagement
$\kappa_q$ initially rises to $0.20$ at $k = 5$ (identity not yet
meaningfully removed), holds for $k \leq 10$, then collapses to
$0.08$ at $k \geq 50$. The signal that LR exploits is concentrated
in the top 5--15 directions of the identity-aligned basis.
\item \textbf{Linear erasure of identity destroys engagement.}
LEACE~\cite{belrose2023leace} applied to SigLIP-L drops subject-ID
accuracy to $0.9\%$ and engagement $\kappa_q$ to $0.04$ in tandem.
Soft erasure (the $k$-direction sweep) inherits the same
trade-off.
\end{itemize}

\paragraph{Interpretation.} Engagement signal in frozen
vision-language features does not occupy a subspace that is
linearly disjoint from identity. It overlaps with the top-rank
identity directions strongly enough that any linear identity-erasure
also erases engagement. Conversely, leaving these directions
intact means the LR readout is fundamentally a subject-prior
classifier with a small engagement perturbation.

\subsection{Mechanism: within-subject ranking is dead}
\label{sec:within}

The identity dominance has a striking behavioural consequence. We
take the predictions of our SOTA recipe (Sec.~\ref{sec:recipe},
$\kappa_q = 0.238$ globally) and re-evaluate them within each test
subject:

\begin{itemize}\itemsep0pt
  \item \textbf{Global $\kappa_q$:} $0.238$
  \item \textbf{Global Spearman $\rho$:} $0.249$
  \item \textbf{Within-subject Spearman $\rho$ (mean over 21 test
  subjects, $\geq 3$ distinct labels):} $0.052$
  \item \textbf{Within-subject AUC ($L = 0$ vs.\ $\geq 1$):} $0.799$
\end{itemize}

In words: the model can tell which student is typically more
engaged than another, and can sometimes tell when a single
student is completely disengaged, but it cannot rank that
student's own clips in order of their engagement. This is the
behavioural signature of a model that has learned a subject-level
prior, not a fine-grained per-clip representation. It is
consistent with the entanglement curve and the linear-erasure
failure.

\paragraph{Implication.} The flat ceiling at $\kappa_q \approx 0.24$
is structural. Readout-level adaptations cannot break it because
the readout is bottlenecked by the encoder's representation, which
has learnt to use identity directions for any informative
discrimination it can make. Either (a)~the encoder must be
fine-tuned to expose identity-orthogonal engagement directions,
or (b)~a non-frozen-feature signal (audio, multi-frame video
temporal, behavioural-explicit MediaPipe signals) must be added.
The remainder of the paper investigates path (a) at the level of
a lightweight intervention---\textsc{Dream}---and reports negative
findings that confirm the structural diagnosis.
